\title{Arduino Anfängerkurs - Abschlussprojekt}
\author{FabLab Ansbach e.V.}
\date{01. April 2026}

\input{../include/template}

\usepackage[T1]{fontenc}
\usepackage[utf8]{inputenc}
\usepackage[ngerman]{babel}
\usepackage{lmodern}
\usepackage{tikz}
\usepackage{xcolor}
\usepackage{mdframed}
\usepackage{array}
\usepackage{tabularx}
\usepackage{parskip}
\usepackage{amssymb}   % \square, \blacktriangleright

% Arduino-Farben
\definecolor{arduinoBlue}{HTML}{00979D}
\definecolor{arduinoDark}{HTML}{005C5F}
\definecolor{arduinoLight}{HTML}{E8F8F8}
\definecolor{borderGray}{HTML}{CCCCCC}
\definecolor{labelColor}{HTML}{444444}

\setlength{\parindent}{0pt}

% Abschnitt-Box (farbiger Rahmen + Titel)
\newmdenv[
  linecolor=arduinoBlue,
  linewidth=1.2pt,
  roundcorner=5pt,
  backgroundcolor=white,
  innertopmargin=7pt,
  innerbottommargin=9pt,
  innerleftmargin=10pt,
  innerrightmargin=10pt,
  leftmargin=6pt,
  rightmargin=6pt,
  skipabove=5pt,
  skipbelow=0pt
]{sectionframe}

% Abschnitt-Box (farbiger Rahmen + Titel)
\newmdenv[
  linecolor=arduinoBlue,
  linewidth=1.2pt,
  roundcorner=5pt,
  backgroundcolor=white,
  innertopmargin=7pt,
  innerbottommargin=9pt,
  innerleftmargin=10pt,
  innerrightmargin=10pt,
  skipabove=5pt,
  skipbelow=0pt
]{sectiontableframe}

% Kopfzeilen-Box (hellblauer Hintergrund)
\newmdenv[
  linecolor=arduinoBlue,
  linewidth=1.2pt,
  roundcorner=5pt,
  backgroundcolor=arduinoLight,
  innertopmargin=7pt,
  innerbottommargin=9pt,
  innerleftmargin=10pt,
  innerrightmargin=10pt,
  skipabove=0pt,
  skipbelow=0pt
]{headerframe}

% Einzelne Schreiblinie
\newcommand{\writeline}{%
  \vspace{6pt}%
  \noindent\tikz\draw[arduinoBlue!45, line width=0.55pt] (0,0) -- (\linewidth,0);%
  \vspace{1pt}%
}

% Nummerierte Schreiblinie
\newcommand{\numberedline}[1]{%
  \vspace{5pt}%
  \noindent\textcolor{arduinoBlue}{\small\textbf{#1.}}\,%
  \tikz\draw[arduinoBlue!45, line width=0.55pt] (0,0) -- (\linewidth-1.1cm,0);%
  \vspace{1pt}%
}

% Checkbox-Schreiblinie
\newcommand{\checkline}{%
  \vspace{5pt}%
  \noindent\textcolor{arduinoBlue}{$\square$}\,%
  \tikz\draw[arduinoBlue!45, line width=0.55pt] (0,0) -- (\linewidth-0.9cm,0);%
  \vspace{1pt}%
}

% Abschnittstitel
\newcommand{\sectiontitle}[2]{%
  {\color{arduinoDark}\textbf{\large #1\enspace #2}}\\[5pt]%
}

\begin{document}

%% NAME & PROJEKTNAME (zweispaltig)
\begin{tabularx}{\linewidth}{X@{\hspace{8pt}}X}
  \begin{headerframe}
    \sectiontitle{$\star$}{Mein Name}
    \writeline\\
    \writeline
  \end{headerframe}
  &
  \begin{headerframe}
    \sectiontitle{$\star$}{Mein Projektname}
    \writeline\\
    \writeline
  \end{headerframe}
\end{tabularx}

\vspace{8pt}

%% PROJEKTBESCHREIBUNG
\begin{sectionframe}
  \sectiontitle{$\blacktriangleright$}{Was macht mein Projekt? \normalfont\small\color{labelColor}(1\,--\,2 Sätze)}
  \writeline\\
  \writeline\\
  \writeline
\end{sectionframe}

\vspace{2pt}

%% SKIZZE
\begin{sectionframe}
  \sectiontitle{$\blacktriangleright$}{Skizze meiner Schaltung}
  \vspace{2pt}
  \begin{center}
    \begin{tikzpicture}
      \draw[borderGray, line width=0.8pt, dashed, rounded corners=4pt]
        (0,0) rectangle (\linewidth-0.4cm, 8cm);
      \node[text=borderGray, font=\small\itshape]
        at ({(\linewidth-0.4cm)/2}, 4cm)
        {Bauteile, Kabel und Arduino hier einzeichnen \ldots};
    \end{tikzpicture}
  \end{center}
\end{sectionframe}

\vspace{2pt}

%% ERSTE SCHRITTE & PROBLEME (zweispaltig)
\begin{tabularx}{\linewidth}{X@{\hspace{8pt}}X}
  \begin{sectiontableframe}
    \sectiontitle{$\blacktriangleright$}{Wie fange ich an?}
    {\small\color{labelColor}Meine ersten Schritte:}
    \newline\\
    \numberedline{1}
    \numberedline{2}
    \numberedline{3}
    \numberedline{4}
    \numberedline{5}
  \end{sectiontableframe}
  &
  \begin{sectiontableframe}
    \sectiontitle{$\blacktriangleright$}{Mögliche Probleme}
    {\small\color{labelColor}Was könnte schwierig werden?}
    \newline\\
    \checkline\\
    \checkline\\
    \checkline\\
    \checkline\\
    \checkline\\
  \end{sectiontableframe}
\end{tabularx}

\vfill
%% FUSSZEILE
\begin{center}
  \textcolor{arduinoBlue!60}{%
    \small\rule{0.25\linewidth}{0.4pt}\quad
    Viel Spa\ss{} beim T\"{u}fteln und Ausprobieren!
    \quad\rule{0.25\linewidth}{0.4pt}%
  }
\end{center}

\end{document}